\documentclass[11pt,a4paper]{article}
\usepackage[utf8]{inputenc}
\usepackage[T1]{fontenc}
\usepackage{lmodern}
\usepackage{textcomp}
\usepackage{amsmath,amssymb,amsthm}
\usepackage{booktabs,array}
\usepackage{enumitem}
\usepackage{xcolor}
\usepackage[expansion=false]{microtype}
\usepackage{url}
\usepackage[margin=2.6cm]{geometry}
\usepackage[pdftitle={Four-colour Rado numbers of a(x-y)=bz},pdfauthor={Takumi Ichikawa},colorlinks=true,linkcolor=blue!50!black,citecolor=blue!50!black,urlcolor=blue!50!black]{hyperref}

\newtheorem{theorem}{Theorem}[section]
\newtheorem{lemma}[theorem]{Lemma}
\newtheorem{proposition}[theorem]{Proposition}
\newtheorem{corollary}[theorem]{Corollary}
\newtheorem{conjecture}[theorem]{Conjecture}
\theoremstyle{definition}
\newtheorem{definition}[theorem]{Definition}
\theoremstyle{remark}
\newtheorem{remark}[theorem]{Remark}

\newcommand{\R}{\mathrm{R}}
\newcommand{\todo}[1]{\textcolor{red}{[#1]}}
\setlist{itemsep=2pt,topsep=4pt}
\allowdisplaybreaks

\title{Four-colour Rado numbers of $a(x-y)=bz$:\\[4pt]\large $\R_4\bigl(3(x-y)=4z\bigr)=461$, $\ \R_4\bigl(5(x-y)=4z\bigr)=741$,\\ $\R_4\bigl(6(x-y)=5z\bigr)=1501$}
\author{Takumi Ichikawa\\ \small\texttt{github.com/IchiTaku326}}
\date{Version 1, 3 November 2026 \\ \small (preprint, not peer reviewed)}
\begin{document}
\maketitle
\begin{abstract}
Let $\R_k(a,b)$ be the least $n$ such that every $k$-colouring of $\{1,\dots,n\}$ has a monochromatic solution of $a(x-y)=bz$ ($\gcd(a,b)=1$).
We determine three new four-colour values, $\R_4(3,4)=461$, $\R_4(5,4)=741$ and $\R_4(6,5)=1501$. The lower bounds are explicit colourings (one of them from a family proved for every number of colours), the upper bounds are SAT computations on a small subset of $[n]$ split by cube-and-conquer, with every proof checked by the formally verified checker \texttt{cake\_lpr}. We also prove $\R_k(a(x-y)=(a-1)z)\ge a^k+a^{k-1}-2a+1$ for all $k,a\ge2$, conjecture equality for $k=4$ (true for $a=3,4,5,6$; for $a=4$ only by a SAT solver without proof file), and collect lower bounds and evidence on a conjectured recursion $\R_4=b(\R_3+2a)+1$ for $a<b$.
\end{abstract}

\section{Introduction}
For positive integers $a,b,k$ the \emph{Rado number} $\R_k(a,b)$ is the least $n$ such that every $k$-colouring of $[n]=\{1,\dots,n\}$ contains a monochromatic solution of $a(x-y)=bz$ (the numbers $x,y,z$ need not be distinct). We always assume $\gcd(a,b)=1$; then every solution has the form
\[
  z=as,\qquad x-y=bs\qquad(s\ge1).
\]
Known four-colour values include $\R_4(1,1)=45$ (one more than the Schur number $S(4)=44$), $\R_4(2,3)=226$ \cite{Ich26a}, and the values listed in \cite{CDW22}. For three colours see \cite{CDW22,AZB25,Ich26b}.

\begin{theorem}\label{thm:main}
\[
  \R_4\bigl(3(x-y)=4z\bigr)=461,\qquad \R_4\bigl(5(x-y)=4z\bigr)=741,\qquad \R_4\bigl(6(x-y)=5z\bigr)=1501 .
\]
\end{theorem}
For each value, the lower bound is an explicit colouring of $[n-1]$ (Section~\ref{sec:lower}) and the upper bound is a computer proof that every colouring of $[n]$ fails (Section~\ref{sec:method}).

\clearpage
\section{Upper bounds: method}\label{sec:method}
\begin{lemma}[small set]\label{lem:small}
Let $K\ge1$ and
\[
  S=\{\,am\le n:\ m\le K\,\}\ \cup\ \{\,x\le n:\ x\equiv n\pmod b\,\}.
\]
If $S$ has no $k$-colouring without monochromatic solutions inside $S$, then $\R_k\le n$. Moreover we may assume that the $i$-th smallest element of $S$ uses one of the colours $0,\dots,i-1$ (``first fit'').
\end{lemma}
\begin{proof}
A good colouring of $[n]$ restricts to $S$. Renaming the colours in order of first appearance gives a first-fit colouring.
\end{proof}

We take $K=\R_3$. The formula for $S$ is about $25$ times smaller than for $[n]$, but still too hard for a single solver run. We therefore use cube-and-conquer \cite{HKM16}:
\begin{enumerate}
\item The \emph{root cubes} are all first-fit colourings of the first $7$ elements of $S$.
\item A cube that is not solved within a time limit is split further by the next $3$ elements of $S$.
\item Every leaf cube is shown unsatisfiable by CaDiCaL, which writes an LRAT proof for the formula with the cube added as unit clauses.
\item Every proof is checked by \texttt{lrat-check} and by the formally verified checker \texttt{cake\_lpr} \cite{THM21}.
\item A separate program checks that the tree is complete, i.e.\ that the leaves cover all first-fit colourings.
\end{enumerate}
In addition, every clause of the encoding is re-checked against the equation by an independent routine.

\begin{center}
\renewcommand{\arraystretch}{1.15}
\begin{tabular}{lrrrl}
\toprule
equation & $n$ & $K=\R_3$ & leaf cubes & checked by \texttt{cake\_lpr}\\
\midrule
$3(x-y)=4z$ & $461$ & $109$ & $22863$ & all\\
$5(x-y)=4z$ & $741$ & $141$ & $813$ & all\\
$6(x-y)=5z$ & $1501$ & $241$ & $6833$ & all\\
\bottomrule
\end{tabular}
\end{center}
Here $\R_3(3,4)=109=b(b^2+b+2a+1)+1$ \cite{AZB25,Ich26b}, and $\R_3(5,4)=141$, $\R_3(6,5)=241$ are the values $a^3+(a-1)^2$ of \cite{CDW22}.

\clearpage
\section{Lower bounds}\label{sec:lower}
\paragraph{The value $461$.} The recursive colouring of \cite[Lemma~4.1 and Theorem~4.5]{Ich26b} for $(a,b)=(3,4)$ and $k=4$ colours $[460]$ without monochromatic solutions. We also checked this colouring against all solutions.

\paragraph{The values $741$ and $1501$.} They are the case $k=4$, $a=5,6$ of the following theorem.

\begin{theorem}\label{thm:H3}
For all $k\ge2$ and $a\ge2$:\quad $\R_k\bigl(a(x-y)=(a-1)z\bigr)\ge a^k+a^{k-1}-2a+1$.
\end{theorem}

\begin{proof}
Let $n=a^k+a^{k-1}-2a$, and let $v(x)$ be the exponent of $a$ in $x$. Since $n<a^{k+1}$ we have $v(x)\le k$.

\emph{The colouring.} A number with $v(x)=j\ge1$ gets colour $c_j$. For the numbers with $v(x)=0$, cut $[1,n]$ into consecutive intervals
\[
  I_2,\ I_3,\ \dots,\ I_{k-1},\ I_k,\ I'_{k-1},\ \dots,\ I'_3,\ I'_2,\qquad |I_j|=|I'_j|=(a-1)a^{j-1},
\]
and give the numbers with $v(x)=0$ in $I_j\cup I'_j$ the colour $c_j$. The total length is $2(a^{k-1}-a)+(a-1)a^{k-1}=n$, and explicitly
\[
  I_j=[\,a^{j-1}-a+1,\ a^j-a\,],\qquad I'_j=[\,n-a^j+a+1,\ n-a^{j-1}+a\,].
\]
Note that colour $c_1$ is used only for numbers with $v(x)=1$.

\emph{No monochromatic solution.} Let $(x,y,z)$ be a solution of colour $c_j$, so $z=as$ and $x-y=(a-1)s$. Since $\gcd(a,a-1)=1$,
\[
  v(z)=1+v(s),\qquad v(x-y)=v(s).
\]
As $v(z)\ge1$, the colour of $z$ is $c_{v(z)}$, so $v(z)=j$ and $v(x-y)=j-1$. Each of $x,y$ has $v=j$, or $v=0$ and $j\ge2$.
\begin{itemize}
\item If $v(x)=v(y)=j$, then $v(x-y)\ge j$, a contradiction.
\item If exactly one of $x,y$ has $v=0$, then $v(x-y)=0$, so $j=1$; but colour $c_1$ is not used for $v=0$.
\item So $v(x)=v(y)=0$, $j\ge2$, and $x-y=(a-1)a^{j-1}u$ with $a\nmid u$ and $z=a^ju\le n$.
  If $x$ and $y$ lie in the same interval of colour $c_j$, then $x-y<(a-1)a^{j-1}$, a contradiction.
  Otherwise $j\le k-1$, $y\in I_j$, $x\in I'_j$, and
  \[
    x-y\ \ge\ (n-a^j+a+1)-(a^j-a)\ =\ a^k+a^{k-1}-2a^j+1 .
  \]
  \begin{itemize}
  \item If $j\le k-2$: $x-y=\frac{a-1}{a}\,z\le\frac{a-1}{a}\,n<a^k-a^{k-2}\le a^k+a^{k-1}-2a^j+1$, a contradiction.
  \item If $j=k-1$: from $a^{k-1}u=z\le n<a^{k-1}(a+1)$ we get $u\le a$, and $a\nmid u$ gives $u\le a-1$. Then $x-y\le(a-1)^2a^{k-2}<a^k-a^{k-1}+1$, a contradiction. \qedhere
  \end{itemize}
\end{itemize}
\end{proof}
For $k=3$ this construction is \cite[Lemma~4.2]{CDW22}, where equality is proved. We also checked the colouring against all solutions for $2\le a\le7$, $2\le k\le5$ (when $n\le20000$).

\begin{conjecture}[H3]
$\R_4\bigl(a(x-y)=(a-1)z\bigr)=a^4+a^3-2a+1$.
\end{conjecture}
This holds for $a=3$ ($103$, known), $a=4$ ($313$, by a SAT solver without proof file) and $a=5,6$ (Theorem~\ref{thm:main}). For $a=7$ the predicted value is $2731$.

\clearpage
\section{The case \texorpdfstring{$a<b$}{a<b}: a recursion and its limits}
In \cite{Ich26b} a recursive colouring gives $\R_k\ge L_k$ with $L_2=b^2+b+1$, $L_k=b(L_{k-1}+2a)+1$ for every $k$, when $D=a^2+ab-b^2$ satisfies $D\le1$ or $a+1\le D\le 2a$ (and $a^2+(b-1)a\le b^2+b+1$). For three colours $L_3$ is the conjectured value of $\R_3$.
\begin{conjecture}[K4]
In the cases above, $\R_4(a,b)=b(\R_3(a,b)+2a)+1$.
\end{conjecture}
Evidence and limits:
\begin{itemize}
\item Equality holds for $(3,4)$ (Theorem~\ref{thm:main}). The recursive colouring gives the further lower bounds
\begin{center}
\begin{tabular}{ccccc}
\toprule
$\R_4(3,5)$ & $\R_4(4,7)$ & $\R_4(5,8)$ & $\R_5(3,4)$ & $\R_5(3,5)$\\
\midrule
$\ge961$ & $\ge3249$ & $\ge5401$ & $\ge1869$ & $\ge4836$\\
\bottomrule
\end{tabular}
\end{center}
all of which are the values $L_4$, $L_5$ predicted by the recursion.
\item If the multiples of $a$ are fixed to the recursive pattern, the remaining constraints on the residue class $x\equiv n\pmod b$ become unsatisfiable exactly at $n=L_4$ for $(a,b)=(3,4)$, $(3,5)$, $(4,7)$, $(5,8)$, $(3,7)$ and $(4,9)$.
\item For $(3,5)$ Lemma~\ref{lem:small} does not decide $961$: $S$ is colourable even with all multiples of $3$ included, because in $S$ the fourth colour can enter the copy $3\cdot[1,186]$ (from $m=170$ in the colouring found). That colouring does not extend to $[961]$, and all our structured searches for a colouring of $[961]$ fail. A further cube-and-conquer run on a reduced set of $479$ elements did not finish either: after splitting, every cube was either refuted within a fraction of a second or not decided within $300$ to $1800$ seconds. We expect $\R_4(3,5)=961$ but have no proof; the upper bound is open.
\item Outside the region the recursion fails: a colouring of $[960]$ for $(4,5)$ (found by SAT with the colour depending only on the $a$-adic valuation and $\lfloor x/2\rfloor$, checked against all solutions) gives $\R_4(4,5)\ge961>941$. In this colouring the fourth colour appears in the copy from $m=156$ on.
\item For $a=2$ the values deviate already for three colours ($\R_3(2,3)=61$, $\R_3(2,5)=181$), and $\R_4(2,3)=226\ne196$; for $(2,5)$ an explicit colouring of $[950]$ (found by SAT, checked against all solutions) gives $\R_4(2,5)\ge951$, well above the value $926$ predicted by the recursion.
\end{itemize}
Five colours: the method of Lemma~\ref{lem:small} did not finish for $\R_5(3,4)\le1869$ (thousands of hard cubes).

\clearpage
\section{Computing environment}
All computations of this note ran on one desktop PC, except the search for the colouring of $[950]$ for $(2,5)$ and some other colourings of Section~4, which ran in a Linux cloud sandbox.
\begin{center}
\small
\renewcommand{\arraystretch}{1.12}
\begin{tabular}{>{\raggedright}p{3.2cm}>{\raggedright\arraybackslash}p{10.4cm}}
\toprule
CPU & Intel Core i5-11400F (6 cores, 12 threads)\\
Memory & 32\,GB DDR4\\
GPU & NVIDIA GeForce RTX 3060 (not used; the SAT solvers run on the CPU)\\
System & Windows~11 Home; solvers and proof checkers in WSL (Linux)\\
Workers & $11$ processes for the cube search and for the proof check\\
Cube search & MiniSat~2.2 through PySAT; cubes it could not solve quickly were tried by CaDiCaL~1.5.3 (PySAT) before splitting\\
Proofs & CaDiCaL~3.0.1 (static build, commit \texttt{c607304}), LRAT\\
Proof checkers & \texttt{lrat-check} (drat-trim, commit \texttt{2e3b2dc}) and \texttt{cake\_lpr} (commit \texttt{a36874a})\\
\bottomrule
\end{tabular}
\end{center}
\begin{center}
\small
\renewcommand{\arraystretch}{1.12}
\begin{tabular}{lrrr}
\toprule
 & $461$ & $741$ & $1501$\\
\midrule
clauses per leaf formula & $31{,}151$--$31{,}160$ & $72{,}288$--$72{,}291$ & $191{,}958$--$191{,}964$\\
leaf cubes, all UNSAT & $22{,}863$ & $813$ & $6{,}833$\\
proofs accepted by \texttt{cake\_lpr} & $22{,}863$ & $813$ & $6{,}833$\\
total size of the LRAT proofs & $196$\,GB & $24.5$\,GB & $256$\,GB\\
cube search, first to last record & $12$\,h\,$56$\,min & $34$\,min & $3$\,h\,$9$\,min\\
proof check, task time (sum over workers) & $34$\,h & $2$\,h & $36$\,h\\
\bottomrule
\end{tabular}
\end{center}
The runs were stopped and resumed several times, so the times in the last two rows are approximate: ``first to last record'' includes pauses, and ``task time'' is the sum of the times of the single tasks over all workers. The proofs were written to a temporary folder and deleted after checking; they can be regenerated with the published scripts.

\section{Reproducibility}
Encodings, cube trees, results of the proof checks, checkers and colourings are available at \url{https://github.com/IchiTaku326/rado-ab}. The LRAT proofs themselves are not stored (they are large); they can be regenerated with the published scripts and checked again. The checker \texttt{cake\_lpr} was compiled with heaps of $4$ and $14$\,GB.

\section*{Acknowledgements and disclosure}
Programs, computation plans, parts of the proofs and parts of the text were developed with the help of an AI assistant (Claude, Anthropic). The author is responsible for the content.

\clearpage
\begin{thebibliography}{99}
\bibitem{AZB25} T.~Ahmed, L.~Zaman, C.~Bright, Symbolic sets for proving bounds on Rado numbers, in: Proc.\ SC-Square 2025, CEUR Workshop Proceedings 4116 (2025); arXiv:2505.12085.
\bibitem{CDW22} Y.~Chang, J.\,A.~De Loera, W.\,J.~Wesley, Rado numbers and SAT computations, arXiv:2210.03262 (2022).
\bibitem{HKM16} M.\,J.\,H.~Heule, O.~Kullmann, V.\,W.~Marek, Solving and verifying the Boolean Pythagorean triples problem via cube-and-conquer, \emph{SAT 2016}.
\bibitem{Ich26a} T.~Ichikawa, The four-colour Rado number of $2(x-y)=3z$ is $226$: a cube-and-conquer computation with checked DRAT proofs, preprint, Zenodo (2026), \url{https://doi.org/10.5281/zenodo.23582301}.
\bibitem{Ich26b} T.~Ichikawa, Rado numbers of $a(x-y)=bz$ with $a<b$: two colours, lower bounds for every number of colours, and the three-colour upper bound up to one statement, companion preprint (2026), \url{https://github.com/IchiTaku326/rado-ab}.
\bibitem{THM21} Y.\,K.~Tan, M.\,J.\,H.~Heule, M.\,O.~Myreen, cake\_lpr: verified propagation redundancy checking in CakeML, in: \emph{TACAS 2021}, LNCS 12652, 223--241.
\end{thebibliography}
\end{document}
